\documentclass{article}
\usepackage[utf8]{inputenc}
\usepackage{amsmath}
\usepackage{amssymb}
\usepackage{amsthm}
\usepackage{algorithm}
\usepackage{algpseudocode}
\usepackage{graphicx}
\usepackage{geometry}
\usepackage{booktabs}
\geometry{a4paper, margin=1in}

% hyperref must be loaded (nearly) last. It subsumes `url` and provides
% the \href command that the `plainurl` bibstyle emits in the .bbl for
% DOIs.
\usepackage[hidelinks]{hyperref}

\newcommand{\old}[1]{{}}
\newtheorem{theorem}{Theorem}
\newtheorem{corollary}{Corollary}
\newtheorem{conjecture}[theorem]{Conjecture}
\newtheorem{lemma}[theorem]{Lemma}
\newtheorem{claim}[theorem]{Claim}
\newtheorem{observation}[theorem]{Observation}
\newtheorem{definition}[theorem]{Definition}

\DeclareMathOperator{\wt}{wt}
\DeclareMathOperator{\MST}{MST}
\DeclareMathOperator{\APSP}{APSP}

\title{Descending Greedy Filter:\\Light Spanners via Edge-Minimality}
\author{P.~Carmi, N.~Horovitz, Y.~Yashuk}
\date{April 2026}

\begin{document}

\maketitle

\begin{abstract}
We present a new construction of geometric spanners, which we call the
\emph{Descending Greedy Filter} (DGF). Starting from the complete graph,
DGF examines edges in \emph{non-ascending} order of length and discards
every edge whose removal does not violate the $t$-spanner condition.
We give the algorithm, its time and space complexity, and conjecture
that its output has $O(n)$ edges and total weight $O(\wt(\MST(P)))$.
We then extract from DGF a general filtering procedure that turns any
$t$-spanner into a \emph{minimal} $t$-spanner (one in which every edge
is essential), and conjecture the same sparsity and weight bounds for
every minimal $t$-spanner in the plane. A recursive binary-search
implementation reduces the number of all-pairs shortest path (APSP)
computations from $\Theta(m)$ to $O(k\log m)$, where $m$ is the input
size and $k$ is the number of edges in the output, giving a near-linear
speedup in the number of APSP calls. Experiments on random point sets
in the unit square with $n\in\{1000,\,2000,\,5000\}$ and
$t\in\{1.001,\,1.01,\,1.1,\,1.5\}$ show that DGF and its variants
produce spanners whose edge count matches the classical greedy spanner
to within a fraction of a percent while using $4$--$7\%$ less total
weight, with the exact stretch preserved.
\end{abstract}

\section{Introduction}
Geometric spanners are a central object in computational geometry. Given
a set $P$ of $n$ points in the plane and a stretch parameter $t\ge 1$, a
geometric $t$-spanner is a graph $G=(P,E)$ in which the shortest-path
distance between any pair of points is at most $t$ times their Euclidean
distance. Such structures have been extensively studied because of their
applications in network design, approximation algorithms, and distributed
systems.

Among the many known constructions, the \emph{greedy spanner} of
Alth\"ofer, Das, Dobkin, Joseph and Soares is widely regarded as the
gold standard: it achieves $O(n)$ edges, total weight
$O(\wt(\MST(P)))$, stretch exactly $t$, and, by a line of work
culminating in Filtser and Solomon, is existentially optimal in the
sense that its sparsity and weight cannot be improved in the worst
case. Empirical studies consistently place the greedy spanner ahead of
Yao graphs, $\Theta$-graphs and WSPD-based constructions on every
quality metric except construction time.

This raises a natural question: is the greedy spanner optimal not only
existentially, but also \emph{instance by instance}? Equivalently, on a
given input, can one find a $t$-spanner with lower total weight (or fewer
edges) than the greedy spanner, without sacrificing the
stretch guarantee? The greedy spanner is \emph{edge-minimal} --- no
single edge can be removed while preserving the stretch --- but being
edge-minimal is a much weaker property than being globally optimal, and
nothing in the greedy construction is tailored to exploit long,
redundant edges created late in the process.

\paragraph{Our contributions.}
We introduce the \emph{Descending Greedy Filter} (DGF), a construction
dual to the greedy spanner, and investigate the class of edge-minimal
$t$-spanners in the plane. Specifically:
\begin{itemize}
  \item \textbf{A new construction.} DGF (Algorithm~\ref{alg:dgf})
    starts from the complete graph and removes edges in
    \emph{non-ascending} order of length whenever doing so preserves
    the stretch $t$. The output is by construction an edge-minimal
    $t$-spanner.
  \item \textbf{A filtering procedure.} We extract from DGF a general
    filter (Algorithm~\ref{alg:dgf-filter}) that converts any
    $t$-spanner into an edge-minimal one, and prove
    (Theorem~\ref{thm:filter-minimal}) that the result is a subset of
    the input and minimal in a pair-witness sense.
  \item \textbf{Conjectures on minimality.} We state four conjectures
    (Conjectures~\ref{conj:dgf-edges}--\ref{conj:min-weight}): DGF
    outputs $O(n)$ edges and $O(\wt(\MST(P)))$ weight, and (more
    strongly) \emph{every} edge-minimal $t$-spanner in the plane has
    these same bounds. We support these conjectures experimentally and
    discuss why the weight conjecture for arbitrary minimal spanners
    is more delicate than the sparsity conjecture.
  \item \textbf{A faster implementation.} We show that the $\Theta(m)$
    APSP calls of the naive DGF can be reduced to $O(k\log m)$, where
    $m=\binom{n}{2}$ is the input size and $k$ is the output size
    (Algorithms~\ref{alg:dgf-bisect-main}--\ref{alg:bisect}). The
    technique is a recursive binary search that certifies contiguous
    batches of removable edges with a single APSP call.
  \item \textbf{Empirical comparison.} On random point sets in the
    unit square, with $n\in\{1000,\,2000,\,5000\}$ and
    $t\in\{1.001,\,1.01,\,1.1,\,1.5\}$, DGF and its variants produce
    spanners with edge counts that match the greedy spanner to
    within a fraction of a percent, but with total weight that is
    consistently $4$--$7\%$ lower, while preserving the exact stretch
    guarantee. We also examine the role of preprocessing: running DGF
    on top of a Yao graph leaves the output essentially unchanged, and
    yields a large runtime speedup for moderate $t$, but loses its
    advantage for very strict $t$ (because Yao itself is then no
    longer sparse), whereas the $\sqrt{t}$-Greedy preprocessor remains
    effective in both regimes.
\end{itemize}

\paragraph{Technical overview.}
The key observation behind DGF is that removing the \emph{longest}
redundant edge first allows shorter edges to inherit its role: a long
edge $(p,q)$ that can be detoured through intermediate points
$p\!\to\!x\!\to\!q$ at stretch $\le t$ is almost always detoured
through edges that are \emph{themselves} part of shorter redundant
detour paths, which DGF will still get to examine later. Greedy, by
contrast, commits to short edges first and can leave long edges
stranded. This ordering heuristic is so aggressive that in most
iterations a large prefix of edges is simultaneously removable; our
binary-search implementation exploits exactly this by bisecting on
prefixes, paying one APSP call per bisection rather than per edge.

For the general minimality question we treat the result of DGF as an
instance of a more general phenomenon: on any input, the set of
edge-minimal $t$-spanners forms a combinatorial family that is poorly
understood. We conjecture that this family always lies inside an
$O(n)$-edge, $O(\wt(\MST))$-weight envelope; a proof of either
conjecture would extend the existential-optimality results of Filtser
and Solomon from the single spanner produced by greedy to the entire
minimality class.

\paragraph{Paper organization.}
Section~\ref{sec:algorithm} defines DGF and analyses its time and space
complexity. Section~\ref{sec:minimal-def} defines minimal
$t$-spanners, and Section~\ref{sec:filtering} shows how DGF can be
used as a filter. Section~\ref{sec:conjectures} collects the four
sparsity and weight conjectures, for DGF and for arbitrary minimal
$t$-spanners, together with a discussion of why the weight conjecture
for arbitrary minimal spanners is delicate.
Section~\ref{sec:binary-search} presents the recursive binary-search
implementation. Section~\ref{sec:experiments} reports experiments,
and Section~\ref{sec:related-work} surveys related work.
Section~\ref{sec:open-problems} concludes with open problems.

\newpage

\section{The Descending Greedy Filter Algorithm}
\label{sec:algorithm}

\algrenewcommand{\algorithmiccomment}[1]{\hspace{3em}// #1}
\begin{algorithm}[h]
\caption{Descending Greedy Filter (DGF)}
\label{alg:dgf}
\begin{algorithmic}[1]
\Statex \textbf{Input:}
\Statex \hspace{1cm} $P$: a set of $n$ points in $\mathbb{R}^d$.
\Statex \hspace{1cm} $t$: a real constant with $t>1$.
\Statex \textbf{Output:}
\Statex \hspace{1cm} A set of edges $E$ such that $(P,E)$ is a $t$-spanner of $P$.

\State Sort the $\binom{n}{2}$ unordered pairs of distinct points in \textbf{non-ascending} order of
       their Euclidean distances (breaking ties arbitrarily) and store them in a list $L$.
\State $E \leftarrow L$
\For{each unordered pair $(p,q) \in L$ in \textbf{non-ascending} order}
    \State $E \leftarrow E \setminus \{(p,q)\}$
    \For{each unordered pair $(r,s) \in \binom{P}{2}$}
        \If{$\delta_{E}(r,s) > t \cdot |rs|$}
            \State $E \leftarrow E \cup \{(p,q)\}$
            \State \textbf{break}
        \EndIf
    \EndFor
\EndFor
\State \Return $E$
\end{algorithmic}
\end{algorithm}

\subsection{Correctness}
At every point during the execution the current edge set $E$ is a
$t$-spanner of $P$: the algorithm starts with the complete graph (which
is trivially a $t$-spanner) and only removes an edge when the resulting
graph is still a $t$-spanner. Hence the returned $E$ is a $t$-spanner.

\subsection{Time complexity}
Let $m=\binom{n}{2}=\Theta(n^2)$ be the initial number of edges, and let
$T_{\APSP}(n,m)$ denote the cost of a single all-pairs shortest
path ($\APSP$) computation on a graph with $n$ vertices and $m$ edges.

Sorting the $m$ edges costs $O(m\log m)=O(n^2 \log n)$.

For every edge we check whether its removal causes a violation of the
$t$-spanner condition. The canonical way to do this is to run $\APSP$
after the tentative removal and compare $\delta_E(r,s)$ with
$t\cdot|rs|$ for every pair. Using Dijkstra from all sources with a
Fibonacci heap, one $\APSP$ call costs
\[
   T_{\APSP}(n,m) \;=\; O\bigl(n\cdot(m+n\log n)\bigr).
\]
Since there are $m$ such checks, the total time is
\[
   O\bigl(m\cdot n\cdot(m+n\log n)\bigr) \;=\; O(n^5)
\]
when $m=\Theta(n^2)$, i.e. the full complete graph is the input.

\subsection{Space complexity}
DGF stores the edge list $L$, the current edge set $E$, and the
data structures needed for APSP (distance matrix, adjacency lists,
priority queue). All of these are $O(m)=O(n^2)$, so the total space
complexity is
\[
   O(n^2).
\]

\newpage

\section{Minimal \texorpdfstring{$t$}{t}-spanners}
\label{sec:minimal-def}

\begin{definition}[minimal $t$-spanner]\label{def:minimal}
Let $P$ be a set of points and $t>1$. An edge set $E\subseteq \binom{P}{2}$
is a \emph{minimal $t$-spanner} of $P$ if
\begin{enumerate}
    \item $(P,E)$ is a $t$-spanner of $P$, and
    \item for every edge $(p,q)\in E$, the set $(E \setminus \{(p,q)\})$ is not a $t$-spanner for $P$. 
\end{enumerate}
Equivalently:
\[
   \forall (p,q)\in E,\; \exists (r,s)\in P\times P\;:\;
   \delta_{E\setminus\{(p,q)\}}(r,s) > t\cdot|rs|.
\]
\end{definition}

Informally, a $t$-spanner is \emph{minimal} (also called \emph{edge-minimal})
when no edge can be removed without violating the stretch guarantee. The
greedy spanner is minimal by construction; most classical constructions
($\Theta$-graphs, Yao graphs, WSPD-based spanners) are not.

\section{Minimizing an arbitrary \texorpdfstring{$t$}{t}-spanner with DGF}
\label{sec:filtering}

The DGF procedure of Algorithm~\ref{alg:dgf} can be applied not only to
the complete graph but to any input edge set that already forms a
$t$-spanner. The only change is that the initial list $L$ contains
only the edges of the input $t$-spanner.

\begin{algorithm}[h]
\caption{Descending Greedy Filter on a $t$-spanner}
\label{alg:dgf-filter}
\begin{algorithmic}[1]
\Statex \textbf{Input:}
\Statex \hspace{1cm} $P$: a set of $n$ points in $\mathbb{R}^d$.
\Statex \hspace{1cm} $t$: a real constant with $t>1$.
\Statex \hspace{1cm} $E_i$: an edge set such that $(P,E_i)$ is a $t$-spanner of $P$.
\Statex \textbf{Output:}
\Statex \hspace{1cm} A set of edges $E$ with $E\subseteq E_i$ such that $(P,E)$ is a minimal $t$-spanner of $P$.

\State Sort $E_i$ in \textbf{non-ascending} order by distance into a list $L$.
\State $E \leftarrow L$
\For{each unordered pair $(p,q) \in L$ in \textbf{non-ascending} order}
    \State $E \leftarrow E \setminus \{(p,q)\}$
    \For{each unordered pair $(r,s) \in \binom{P}{2}$}
        \If{$\delta_{E}(r,s) > t \cdot |rs|$}
            \State $E \leftarrow E \cup \{(p,q)\}$
            \State \textbf{break}
        \EndIf
    \EndFor
\EndFor
\State \Return $E$
\end{algorithmic}
\end{algorithm}

\begin{theorem}\label{thm:filter-minimal}
Let $E_i$ be a $t$-spanner of $P$, and let $E$ be the output of
Algorithm~\ref{alg:dgf-filter} on input $(P,t,E_i)$. Then
\begin{enumerate}
    \item $E \subseteq E_i$,
    \item $(P,E)$ is a $t$-spanner of $P$, and
    \item $E$ is a minimal $t$-spanner of $P$ in the sense of
          Definition~\ref{def:minimal}.
\end{enumerate}
\end{theorem}

\begin{proof}
Property (1) is immediate, since the algorithm only removes edges.

Property (2) follows from the invariant that $E$ is a $t$-spanner at
every iteration: the algorithm starts with the $t$-spanner $E_i$, and
whenever an edge is removed we verify that the result is still a
$t$-spanner; otherwise the edge is restored.

Property (3). Let $(p,q)\in E$ be an edge in the output. We must
exhibit a pair $(r,s)\in P\times P$ such that
$\delta_{E\setminus\{(p,q)\}}(r,s)>t\cdot|rs|$.

Denote by $E'$ the state of $E$ at the moment the main loop inspected
the edge $(p,q)$. Because the algorithm only removes edges after the
inspection of $(p,q)$, we have $E\subseteq E'$.

The fact that $(p,q)\in E$ means that the inner violation check
succeeded, i.e.\ there exist $(r,s)\in P\times P$ with
\[
   \delta_{E'\setminus\{(p,q)\}}(r,s) \;>\; t\cdot|rs|.
\]
Since $E\subseteq E'$, removing more edges can only increase shortest
path distances, hence
\[
   \delta_{E\setminus\{(p,q)\}}(r,s) \;\ge\;
   \delta_{E'\setminus\{(p,q)\}}(r,s) \;>\; t\cdot|rs|,
\]
which is exactly the minimality condition for $(p,q)$.
\end{proof}

\subsection{Complexity of the filtering version}
Let $m_i=|E_i|$ denote the number of edges in the input $t$-spanner.
Sorting costs $O(m_i\log m_i)$. For each of the $m_i$ candidate edges
we run an APSP on the current graph with at most $m_i$ edges, giving
\[
   O\bigl(m_i\cdot n\cdot(m_i+n\log n)\bigr)
\]
total time. In the typical case where the input already satisfies
$m_i=O(n)$ (e.g.\ $E_i$ is a Yao, $\Theta$, or greedy spanner), this
simplifies to $O(n^3 \log n)$. The space complexity is $O(m_i)$.

\section{Sparsity and Weight Conjectures}
\label{sec:conjectures}

We now state our main conjectures. Throughout this section, $P$ is a
set of $n$ points in the plane and $t>1$ is a constant; the hidden
constants in the $O$-notation may depend on $t$. The conjectures come
in two flavors. The first pair concerns the specific output of DGF.
The second pair concerns \emph{any} minimal $t$-spanner in the sense
of Definition~\ref{def:minimal}.

\subsection{Conjectures for DGF}
Let $G_{\mathrm{DGF}}=(P,E_f)$ be the output of Algorithm~\ref{alg:dgf}
on $P$ with parameter $t$.

\begin{conjecture}[sparsity of DGF]\label{conj:dgf-edges}
$|E_f| \;=\; O(n)$.
\end{conjecture}

\begin{conjecture}[weight of DGF]\label{conj:dgf-weight}
$\wt(E_f) \;=\; O\!\bigl(\wt(\MST(P))\bigr)$.
\end{conjecture}

\subsection{Conjectures for arbitrary minimal \texorpdfstring{$t$}{t}-spanners}
Let $E$ be any minimal $t$-spanner of $P$ in the sense of
Definition~\ref{def:minimal}.

\begin{conjecture}[sparsity of any minimal $t$-spanner]\label{conj:min-edges}
$|E| \;=\; O(n)$.
\end{conjecture}

\begin{conjecture}[weight of any minimal $t$-spanner]\label{conj:min-weight}
$\wt(E) \;=\; O\!\bigl(\wt(\MST(P))\bigr)$.
\end{conjecture}

\subsection{A conjecture toward Conjecture~\ref{conj:min-weight}}
\label{subsec:toward-min-weight}
We now state an additional conjecture which, if established, we believe
would imply (or substantially advance a proof of)
Conjecture~\ref{conj:min-weight}.

%% The setup, notation, construction and claim below were transcribed
%% from a whiteboard sketch; please double-check the definitions and
%% the form of $T_{pq}$ before publication.

\paragraph{Setup.}
Let $E_Z = \binom{P}{2}$ denote the edge set of the complete graph on
$P$, and let $E \subseteq E_Z$ be a minimal $t$-spanner of $P$ in the
sense of Definition~\ref{def:minimal}. For a subset $F \subseteq E_Z$
and points $u,v \in P$, write $\delta_F(u,v)$ for the shortest-path
distance from $u$ to $v$ in the weighted graph $(P,F)$ (using Euclidean
edge lengths).

\paragraph{Notation.}
For each edge $(p,q) \in E$, define
\begin{align*}
  C_{pq} &\;=\; \Bigl\{(a,b) \in E \;\Bigm|\;
    \delta_{E\setminus\{(a,b)\}}(a,b) > t\,|ab|,\
    \text{$(p,q)$ is a ``bottleneck'' of $(a,b)$}\Bigr\}, \\
  (a,b)_{pq} &\;=\; \arg\min_{(a,b)\,\in\, C_{pq}} |ab|, \\
  T_{pq} &\;=\; \delta_{E_Z}(a,p) \;+\; |pq| \;+\; \delta_{E_Z}(q,b),
  \qquad \text{where $(a,b)=(a,b)_{pq}$.}
\end{align*}
%% NOTE: ``$(p,q)$ is a bottleneck of $(a,b)$'' was written informally
%% on the board; the intended meaning (to verify) is that every
%% $t$-path between $a$ and $b$ in $E$ must use the edge $(p,q)$.

\paragraph{Construction of $E^*$.}
Initialize $E^* \gets \emptyset$. For every $(p,q) \in E$, add the
shortest path realizing $T_{pq}$ to $E^*$, i.e.\
$E^* \gets E^* \cup T_{pq}$.

\begin{conjecture}[auxiliary conjecture toward weight of any minimal $t$-spanner]
\label{conj:toward-min-weight}
With the construction above,
\[
  \wt(E^*) \;\leq\; 2\,\wt(\MST(P)).
\]
\end{conjecture}

\paragraph{Why this would help.}
%% TODO: write up the implication. Whiteboard hint:
%% ``Proof: for every $T_i, T_j \in E^*$ \ldots\ MST \ldots'',
%% i.e.\ a pairwise / charging argument against the MST.
If Conjecture~\ref{conj:toward-min-weight} holds, then every minimal
$t$-spanner $E$ admits a witness collection $\{T_{pq}\}_{(p,q)\in E}$
whose total weight is at most $2\,\wt(\MST(P))$; combined with the
fact that each edge $(p,q)\in E$ contributes its own length $|pq|$ to
some $T_{pq}$, this would yield
$\wt(E) = O\!\bigl(\wt(\MST(P))\bigr)$ and hence
Conjecture~\ref{conj:min-weight}.

\subsection{Discussion}
Conjectures~\ref{conj:min-edges} and~\ref{conj:min-weight} strictly
generalize Conjectures~\ref{conj:dgf-edges} and~\ref{conj:dgf-weight}
respectively, since the output of DGF is minimal by
Theorem~\ref{thm:filter-minimal} (applied to $E_i=\binom{P}{2}$).
Likewise, the greedy spanner is minimal by construction, so the
conjectures are consistent with the classical bounds $O(n)$ edges
and $O(\wt(\MST(P)))$ weight that are known for the greedy spanner.
All four conjectures are supported by the experiments reported in
Section~\ref{sec:experiments}: in every tested instance the number of
edges of each minimal spanner we produce is close to $n$ and its
total weight is at most a small multiple of $\wt(\MST(P))$.

\newpage

\section{Faster Implementation via Recursive Binary Search}
\label{sec:binary-search}

The na\"ive implementation of DGF (Algorithm~\ref{alg:dgf}) iterates
over all edges one by one, and for each edge performs an APSP
computation. This results in $\Theta(m)$ APSP calls, where
$m=\binom{n}{2}$. In practice, the vast majority of the long edges turn
out to be removable, and only a small subset is necessary. This
motivates a binary-search approach that can certify large batches of
removable edges with a single APSP call.

\subsection{Key observation}
Let $E$ be the current edge set (sorted in non-ascending order by
weight), and let $S \subseteq E$ be a contiguous batch of edges. If
removing all edges of $S$ from $E$ does not violate the $t$-spanner
condition---that is, if for every $r,s\in P$ we have
$\delta_{E\setminus S}(r,s) \le t\cdot|rs|$---then every edge in $S$
is removable. A single APSP call therefore certifies the removal of an
entire batch. When a violation is detected, we split the batch into two
halves and recurse on each.

\subsection{Algorithm}
\begin{algorithm}[h]
\caption{DGF with recursive binary search}
\label{alg:dgf-bisect-main}
\begin{algorithmic}[1]
\Statex \textbf{Input:} $P$, $t$ as in Algorithm~\ref{alg:dgf}.
\Statex \textbf{Output:} a $t$-spanner $E$ of $P$ identical to the output of Algorithm~\ref{alg:dgf}.

\State Sort the $\binom{n}{2}$ pairs in \textbf{non-ascending} order by distance into a list $L$.
\State $E \leftarrow L$
\State $E \leftarrow \textsc{Bisect}(\emptyset,\, L,\, P,\, t)$
\State \Return $E$
\end{algorithmic}
\end{algorithm}

\begin{algorithm}[h]
\caption{\textsc{Bisect}$(E,\,S,\,P,\,t)$}
\label{alg:bisect}
\begin{algorithmic}[1]
\Statex \textbf{Input:}
\Statex \hspace{1cm} $E$: current edge set (i.e.\ $S$ has been removed from $E$ before the call).
\Statex \hspace{1cm} $S$: batch of candidate edges to confirm as removable, sorted in non-ascending order by weight.
\Statex \hspace{1cm} $P$: point set.
\Statex \hspace{1cm} $t$: stretch parameter.
\Statex \textbf{Output:} updated $E$ in which the removable edges of $S$ are discarded and the necessary edges are restored.

\If{$S = \emptyset$}
    \State \Return $E$
\EndIf
\State Compute APSP on $(P,E)$.
\If{$\delta_E(r,s) \le t \cdot |rs|$ for every $(r,s)\in P\times P$}
    \State \Return $E$ \Comment{No violation: every edge in $S$ is removable}
\EndIf
\If{$|S|=1$}
    \State $E \leftarrow E \cup S$ \Comment{Single necessary edge: restore it}
    \State \Return $E$
\EndIf
\State Split $S$ into $S_1$ (first, longer half) and $S_2$ (second, shorter half).
\State $E \leftarrow E \cup S$ \Comment{Restore all edges in $S$}
\State $E \leftarrow E \setminus S_1$
\State $E \leftarrow \textsc{Bisect}(E,\, S_1,\, P,\, t)$
\State $E \leftarrow E \setminus S_2$
\State $E \leftarrow \textsc{Bisect}(E,\, S_2,\, P,\, t)$
\State \Return $E$
\end{algorithmic}
\end{algorithm}

\subsection{Correctness}
The binary-search implementation produces the same output as
Algorithm~\ref{alg:dgf}. Each edge is either certified as removable
(when the entire batch that contains it passes the violation check) or
identified as necessary (when the recursion reaches a single-edge base
case with a violation). Because the edges are processed in the same
non-ascending order and the removal/restoration operations are
symmetric, the final edge set is identical to that of
Algorithm~\ref{alg:dgf}.

\subsection{Time complexity}
Let $m$ denote the number of edges given to \textsc{Bisect} at the top
level, and let $k$ denote the number of necessary edges in the output.
Let $T_{\APSP}(n,m)=O(n(m+n\log n))$ be the cost of a single
$\APSP$ computation.

\paragraph{APSP calls.}
Each necessary edge traces a root-to-leaf path of depth at most
$\lceil\log_2 m\rceil$ in the recursion tree, contributing
$O(\log m)$ APSP calls. Every such path terminates at a single-edge
base case where the edge is restored. Conversely, groups of removable
edges are eliminated at internal nodes without further recursion.
Therefore the total number of APSP calls is
\[
   O(k\log m).
\]

\paragraph{Overall complexity.}
Sorting the edges takes $O(m\log m)$. The binary search performs
$O(k\log m)$ APSP calls, each costing $O(n(m+n\log n))$. In
addition, every level of the recursion performs $O(m)$ work for edge
removal and restoration. The total time is therefore
\[
   O\!\left(m\log m \;+\; k\log m\cdot n(m+n\log n)\right).
\]

\paragraph{Comparison with the na\"ive approach.}
The na\"ive implementation performs $m$ APSP calls, for
$O(m\cdot n(m+n\log n))$. The binary search reduces this to
$O(k\log m\cdot n(m+n\log n))$. Since geometric $t$-spanners are
expected to satisfy $k=O(n)$ while $m=\Theta(n^2)$, the speedup factor
is
\[
   \frac{m}{k\log m} \;=\; \frac{\Theta(n^2)}{O(n\log n)}
   \;=\; \Omega\!\left(\frac{n}{\log n}\right),
\]
a near-linear improvement in the number of APSP calls.

\paragraph{Concrete bounds assuming $k=O(n)$.}
\begin{itemize}
    \item \textbf{Complete graph input} ($m=\binom{n}{2}=\Theta(n^2)$):
    each APSP costs $O(n^3)$, so the total time is
    \[
        O(n^2\log n + n\log n\cdot n^3) \;=\; O(n^4 \log n).
    \]
    \item \textbf{$t$-spanner input} ($m=O(n)$): each APSP costs
    $O(n^2\log n)$, so the total time is
    \[
        O(n\log n + n\log n\cdot n^2\log n) \;=\; O(n^3 \log^2 n).
    \]
\end{itemize}

\newpage

\section{Experiments}
\label{sec:experiments}

\subsection{Setup}
All experiments are run on point sets drawn uniformly at random in the
unit square $[0,1]^2$, with fixed random $\mathtt{seed} (=42)$. We
compare the following algorithms:
\begin{itemize}
    \item \textbf{Greedy}: the classical greedy
          $t$-spanner.
    \item \textbf{$\sqrt{t}$-Greedy + DGF}: run the greedy algorithm
          with parameter $\sqrt{t}$, then apply DGF
          (Algorithm~\ref{alg:dgf-filter}) with parameter $t$ to the
          result.
    \item \textbf{Yao}: the Yao graph with the smallest $k$ such that
          the construction guarantees stretch at most $t$.
    \item \textbf{Yao + DGF}: run Yao with the above $k$, then apply
          DGF with parameter $t$ to the resulting graph.
    \item \textbf{$\sqrt{t}$-Yao + $\sqrt{t}$-Greedy}: run Yao with
          parameter $\sqrt{t}$, then run greedy on the resulting edges
          with parameter $\sqrt{t}$.
    \item \textbf{DGF}: Algorithm~\ref{alg:dgf} starting from the
          complete graph (using the recursive binary-search
          implementation of Section~\ref{sec:binary-search}).
\end{itemize}

For every instance we report the number of edges, the total weight,
the weight ratio over the MST (that is,
$\wt(E)/\wt(\MST(P))$), the maximum and average vertex degrees, the
empirically measured maximum stretch, whether the output is a valid
$t$-spanner, and the running time in seconds.

\subsection{Results for \texorpdfstring{$n=5000,\; t=1.5$}{n=5000, t=1.5}}

\begin{table}[htbp]
    \centering
    \begin{tabular}{lrrrrrrr}
        \toprule
        Algorithm & Edges & Weight & WtRatio & MaxDeg & AvgDeg & MaxStrch & Valid \\
        \midrule
        Yao + DGF                             & 9\,377           & \textbf{128.07}  & \textbf{2.80} & \textbf{7} & 3.75          & 1.5000           & \checkmark \\
        $\sqrt{t}$-Greedy + DGF               & \textbf{9\,314}  & 128.91           & 2.81          & \textbf{7} & \textbf{3.73} & 1.4999           & \checkmark \\
        Greedy                                & 9\,385           & 137.58           & 3.00          & \textbf{7} & 3.75          & 1.4999           & \checkmark \\
        $\sqrt{t}$-Yao + $\sqrt{t}$-Greedy    & 14\,269          & 255.68           & 5.58          & 11         & 5.71          & 1.2754           & \checkmark \\
        Yao                                   & 66\,839          & 2\,302.90        & 50.26         & 52         & 26.74         & \textbf{1.1647}  & \checkmark \\
        \bottomrule
    \end{tabular}
    \caption{Summary for $n=5000$, $t=1.5$, $\mathtt{seed}=42$. Every output is a valid $1.5$-spanner.}
    \label{tab:n5000-t1.5}
\end{table}

Table~\ref{tab:n5000-t1.5} reports the results. For $t=1.5$, both
DGF-based variants produce spanners whose edge counts are essentially
identical to the greedy spanner ($9\,377$ for Yao + DGF and $9\,314$
for $\sqrt{t}$-Greedy + DGF, vs.\ $9\,385$ for greedy---within
$0.1\%$). Their total weight, however, is about $7\%$ lower than
greedy's ($128.07$ and $128.91$ vs.\ $137.58$), and the weight ratio
over the MST drops from $3.00$ to $2.80$.

\subsection{Results for \texorpdfstring{$n=5000,\; t=1.1$}{n=5000, t=1.1}}

\begin{table}[htbp]
    \centering
    \begin{tabular}{lrrrrrrr}
        \toprule
        Algorithm & Edges & Weight & WtRatio & MaxDeg & AvgDeg & MaxStrch & Valid \\
        \midrule
        Yao + DGF                             & 22\,191           & \textbf{463.89}   & \textbf{10.12} & 17           & 8.88          & 1.1000          & \checkmark \\
        $\sqrt{t}$-Greedy + DGF               & \textbf{22\,127}  & 467.67            & 10.21          & \textbf{16}  & \textbf{8.85} & 1.1000          & \checkmark \\
        Greedy                                & 22\,225           & 492.74            & 10.75          & \textbf{16}  & 8.89          & 1.1000          & \checkmark \\
        $\sqrt{t}$-Yao + $\sqrt{t}$-Greedy    & 32\,320           & 858.65            & 18.74          & 25           & 12.93         & 1.0488          & \checkmark \\
        Yao                                   & 224\,442          & 14\,250.23        & 311.00         & 137          & 89.78         & \textbf{1.0378} & \checkmark \\
        \bottomrule
    \end{tabular}
    \caption{Summary for $n=5000$, $t=1.1$, $\mathtt{seed}=42$. Every output is a valid $1.1$-spanner.}
    \label{tab:n5000-t1.1}
\end{table}

Table~\ref{tab:n5000-t1.1} reports the results. For the stricter
stretch $t=1.1$ the pattern is the same: the DGF-based variants and
greedy agree on edge count to within $0.15\%$ ($22\,191$ and
$22\,127$ vs.\ $22\,225$), while the total weight is about $6\%$
lower ($463.89$ and $467.67$ vs.\ $492.74$). The weight ratio over the
MST drops from $10.75$ (greedy) to $10.12$ (Yao + DGF).

\subsection{Results for \texorpdfstring{$n=2000,\; t=1.01$}{n=2000, t=1.01}}

\begin{table}[htbp]
    \centering
    \begin{tabular}{lrrrrrrr}
        \toprule
        Algorithm & Edges & Weight & WtRatio & MaxDeg & AvgDeg & MaxStrch & Valid \\
        \midrule
        DGF                                   & 27\,714           & \textbf{1\,647.42} & \textbf{56.85}  & \textbf{44}  & 27.71          & 1.0100          & \checkmark \\
        Yao + DGF                             & 27\,714           & \textbf{1\,647.42} & \textbf{56.85}  & \textbf{44}  & 27.71          & 1.0100          & \checkmark \\
        $\sqrt{t}$-Greedy + DGF               & \textbf{27\,659}  & 1\,657.07          & 57.18           & \textbf{44}  & \textbf{27.66} & 1.0100          & \checkmark \\
        Greedy                                & 27\,739           & 1\,723.16          & 59.46           & \textbf{44}  & 27.74          & 1.0100          & \checkmark \\
        $\sqrt{t}$-Yao + $\sqrt{t}$-Greedy    & 38\,837           & 2\,864.12          & 98.83           & 57           & 38.84          & 1.0050          & \checkmark \\
        Yao                                   & 571\,029          & 156\,774.77        & 5\,409.63       & 860          & 571.03         & \textbf{1.0008} & \checkmark \\
        \bottomrule
    \end{tabular}
    \caption{Summary for $n=2000$, $t=1.01$, $\mathtt{seed}=42$. Every output is a valid $1.01$-spanner.}
    \label{tab:n2000-t1.01}
\end{table}

Table~\ref{tab:n2000-t1.01} reports the results. For the strict
stretch $t=1.01$, Yao + DGF and the full DGF converge to an
\emph{identical} solution ($27\,714$ edges, weight $1\,647.42$),
despite starting from very different graphs: Yao has $571\,029$ edges
here, roughly $20\times$ the output size, yet every one of the
``extra'' edges is identified as redundant by DGF. The resulting edge
count matches greedy ($27\,739$) to within $0.1\%$, while the total
weight is about $4\%$ lower and the weight ratio over the MST drops
from $59.46$ (greedy) to $56.85$ (DGF / Yao + DGF). The
$\sqrt{t}$-Greedy + DGF variant sits at $27\,659$ edges and weight
$1\,657.07$, very close to the other DGF-based outputs.

\subsection{Results for \texorpdfstring{$n=1000,\; t=1.001$}{n=1000, t=1.001}}

\begin{table}[htbp]
    \centering
    \begin{tabular}{lrrrrrrr}
        \toprule
        Algorithm & Edges & Weight & WtRatio & MaxDeg & AvgDeg & MaxStrch & Valid \\
        \midrule
        DGF                                   & 38\,633          & \textbf{5\,556.17} & \textbf{267.16}  & 127           & 77.27          & 1.0010          & \checkmark \\
        Yao + DGF                             & 38\,633          & \textbf{5\,556.17} & \textbf{267.16}  & 127           & 77.27          & 1.0010          & \checkmark \\
        $\sqrt{t}$-Greedy + DGF               & \textbf{38\,574} & 5\,588.94          & 268.74           & 122           & \textbf{77.15} & 1.0010          & \checkmark \\
        Greedy                                & 38\,682          & 5\,786.30          & 278.23           & \textbf{117}  & 77.36          & 1.0010          & \checkmark \\
        $\sqrt{t}$-Yao + $\sqrt{t}$-Greedy    & 52\,375          & 9\,167.11          & 440.79           & 164           & 104.75         & 1.0005          & \checkmark \\
        Yao                                   & 465\,760         & 235\,769.69        & 11\,336.68       & 982           & 931.52         & \textbf{1.0001} & \checkmark \\
        \bottomrule
    \end{tabular}
    \caption{Summary for $n=1000$, $t=1.001$, $\mathtt{seed}=42$. Every output is a valid $1.001$-spanner.}
    \label{tab:n1000-t1.001}
\end{table}

Table~\ref{tab:n1000-t1.001} reports the results in the most extreme
regime we tested ($t-1=10^{-3}$). The qualitative pattern of the
previous tables is preserved: DGF and Yao + DGF again converge to an
\emph{identical} solution ($38\,633$ edges, weight $5\,556.17$),
matching greedy on edges to within $0.13\%$ ($38\,682$) and reducing
the total weight by about $4\%$ ($5\,786.30 \to 5\,556.17$), with the
weight ratio over the MST dropping from $278.23$ to $267.16$. The
$\sqrt{t}$-Greedy + DGF variant achieves the smallest edge count
($38\,574$, again within $0.3\%$ of greedy).

A new effect appears at this extreme stretch: the Yao graph itself is
no longer a \emph{sparse} preprocessor. With $t-1=10^{-3}$, the cone
half-angle becomes very small, and Yao retains $465\,760$
edges---over $93\%$ of the $\binom{1000}{2}=499\,500$ candidate
pairs---so the input handed to DGF after the Yao step is essentially
the complete graph again. As a consequence, the runtime advantage of
Yao preprocessing observed in the previous regimes disappears here:
running DGF directly is no slower than running Yao + DGF (in our
implementation, both took on the order of $4$--$5\times 10^{4}$
seconds). By contrast, $\sqrt{t}$-Greedy remains effective as a
preprocessor in this regime, because it produces a graph with only
$O(n)$ edges (essentially the greedy $\sqrt{t}$-spanner) and the DGF
pass on top of it terminates in roughly $7\times 10^{3}$
seconds---about an order of magnitude faster than DGF on the full
graph. The output of $\sqrt{t}$-Greedy + DGF is again within a tiny
margin of pure DGF on every quality metric.

\newpage

\subsection{Discussion}

To make the DGF-vs.-greedy comparison precise and uniform across the
four regimes, Table~\ref{tab:summary-percent} reports, for each
metric and each $(n,t)$ pair, the percentage change between the
\emph{best} DGF-based algorithm (over $\{\text{DGF},\,$Yao + DGF$,\,
\sqrt{t}\text{-Greedy} + \text{DGF}\}$) and the \emph{best}
greedy-based algorithm (over $\{$Greedy$,\,\sqrt{t}\text{-Yao} +
\sqrt{t}\text{-Greedy}\}$). Negative numbers favor the DGF family. In
our data the best greedy-based algorithm turns out to be plain Greedy
on every metric in every configuration ($\sqrt{t}$-Yao +
$\sqrt{t}$-Greedy is strictly worse than Greedy across the board), so
each cell is effectively ``best DGF-based vs.\ Greedy'', with the
DGF representative chosen per cell. The rightmost column averages
each row across the four regimes.

\begin{table}[htbp]
    \centering
    \begin{tabular}{lrrrrr}
        \toprule
        Metric  & $t=1.5$ & $t=1.1$ & $t=1.01$ & $t=1.001$ & Average \\
        \midrule
        Edges   & $-0.76\%$           & $-0.44\%$           & $-0.29\%$           & $-0.28\%$ & $-0.44\%$ \\
        Weight  & $-6.91\%$           & $-5.86\%$           & $-4.40\%$           & $-3.98\%$ & $-5.28\%$ \\
        WtRatio & $-6.67\%$           & $-5.86\%$           & $-4.39\%$           & $-3.98\%$ & $-5.22\%$ \\
        MaxDeg  & $\phantom{+}0.00\%$ & $\phantom{+}0.00\%$ & $\phantom{+}0.00\%$ & $+4.27\%$ & $+1.07\%$ \\
        AvgDeg  & $-0.53\%$           & $-0.45\%$           & $-0.29\%$           & $-0.27\%$ & $-0.39\%$ \\
        \bottomrule
    \end{tabular}
    \caption{Percentage change in each metric: best DGF-based output
    minus best greedy-based output, normalized by the best
    greedy-based output. Negative numbers favor the DGF family. The
    Average column is the arithmetic mean across the four regimes.}
    \label{tab:summary-percent}
\end{table}

The table makes the qualitative findings visible at a glance.

\paragraph{Edges (and average degree): ${\le}\,0.76\%$ reduction in every regime.}
On edges, the best DGF-based algorithm in every cell is $\sqrt{t}$-Greedy
+ DGF, and it shaves between $0.28\%$ and $0.76\%$ off greedy's edge
count, for an average reduction of $0.44\%$. Average degree mirrors
this finding (since $\mathrm{AvgDeg} = 2\cdot \mathrm{Edges}/n$),
with an average reduction of $0.39\%$. In absolute terms these are
all changes of well under one percent; for practical purposes we
view edge-minimal $t$-spanners of random Euclidean inputs as having
an edge count that is essentially determined by the input, largely
independent of which edge-minimization strategy produces them. This
is also consistent with what one would expect from
Conjecture~\ref{conj:min-edges}: if \emph{every} edge-minimal
$t$-spanner has $O_{t,d}(n)$ edges, then greedy and DGF should be
close.

\paragraph{Weight (and weight ratio): $4$--$7\%$ reduction, monotone in $t$.}
On total weight, the best DGF-based algorithm is consistently and
measurably lighter than greedy: the gap shrinks monotonically from
$-6.91\%$ at $t=1.5$ to $-3.98\%$ at $t=1.001$, averaging $-5.28\%$
across the four regimes. The weight ratio over the MST drops by
essentially the same amount ($-5.22\%$ on average), since both
quantities are normalized by the same $\wt(\MST)$. The DGF
representative achieving the best weight is Yao + DGF at $t=1.5,1.1$
and pure DGF (which equals Yao + DGF on those instances) at
$t=1.01,1.001$. Descending processing, which removes the longest
redundant edge first, systematically leaves a slightly lighter
edge-minimal representative of the set of all edge-minimal
$t$-spanners on the input. This is the main empirical advantage of
DGF over the classical greedy spanner.

\paragraph{Maximum degree: ${\le}\,4.27\%$ \emph{increase}, only at the strictest $t$.}
The one quality metric on which DGF does \emph{not} match or beat
greedy is the maximum vertex degree, but the gap is small and
isolated. The best DGF-based MaxDeg ties greedy exactly at
$t=1.5,\;1.1$ and $1.01$ (so the comparison is $0.00\%$ in three of
the four regimes), and at $t=1.001$ the best DGF-based MaxDeg
($\sqrt{t}$-Greedy + DGF, with $122$) is $4.27\%$ above greedy's
$117$. The average MaxDeg increase across the four regimes is
$+1.07\%$. We do not currently have a structural explanation for the
pattern; one possible reading is that descending edge processing
tends to leave a small number of slightly-high-degree ``hub'' vertices
through which long-edge detours are routed.

\paragraph{Yao preprocessing: pure speedup for moderate $t$, ineffective for very strict $t$.}
The Yao graph is far from edge-minimal---$66\,839$ edges at $t=1.5$,
$224\,442$ at $t=1.1$, $571\,029$ at $t=1.01$, and $465\,760$ at
$t=1.001$---but it provably contains a valid $t$-spanner, so feeding
it to DGF is safe. Empirically, the DGF outputs obtained from Yao are,
for every $(n,t)$ we tested, either identical to the full-input DGF
output (at $t=1.01$ and $t=1.001$, both produce exactly the same
edge set and weight) or differ by a handful of edges and a negligible
amount of weight. So the preprocessing step is, on output quality,
\emph{essentially neutral}.

On runtime, however, the picture has two regimes. For moderate $t$
($t\ge 1.01$ in our experiments), Yao reduces the number of candidate
edges from $\binom{n}{2}$ to a much smaller subset---e.g.\ from about
$12.5$ million to about $571\,000$ at $n=5000,\;t=1.01$, a $96\%$
reduction---which, via the recursive binary-search implementation
(Section~\ref{sec:binary-search}), yields a large practical speedup
for DGF. For very strict $t$ ($t=1.001$ in our experiments), Yao's
cones become so narrow that it retains essentially the entire
complete graph ($93\%$ of $\binom{1000}{2}$), and the speedup
vanishes: in our implementation Yao + DGF was actually marginally
\emph{slower} than pure DGF in this regime. The $\sqrt{t}$-Greedy
preprocessing remains effective in both regimes, since it always
produces a graph with $O(n)$ edges; at $t=1.001$ it gave roughly an
order-of-magnitude runtime speedup over pure DGF on the same input,
again with essentially the same output quality.

\paragraph{Support for the conjectures.}
The weight ratios over the MST stay bounded and scale with $t$ as
classical theory predicts: $\rho\approx 2.8$ for $t=1.5$,
$\rho\approx 10$ for $t=1.1$, $\rho\approx 57$ for $t=1.01$ and
$\rho\approx 267$ for $t=1.001$, consistent with the $1/(t-1)$
behavior observed for the greedy spanner. The edge counts remain
linear in $n$ across all four $(n,t)$ pairs. Taken together, these
observations give direct empirical support to
Conjectures~\ref{conj:dgf-edges}, \ref{conj:dgf-weight},
\ref{conj:min-edges} and~\ref{conj:min-weight}: every edge-minimal
$t$-spanner we compute---including outputs that are provably
\emph{not} the greedy spanner---has a linear number of edges and
weight $O(\wt(\MST))$.

\paragraph{Bottom line.}
Comparing the best of the DGF family against the classic greedy
spanner on random Euclidean inputs, DGF-based variants reduce total
weight by $5.28\%$ on average (and $3.98$--$6.91\%$ across the four
regimes), reduce edge count by $0.44\%$ on average, and leave the
maximum vertex degree essentially unchanged ($+1.07\%$ on average,
zero in three of four regimes, and at most $+4.27\%$ at the strictest
$t$). The choice of preprocessing affects only runtime, not output
quality: Yao + DGF is a large practical speedup for moderate $t$ and
roughly neutral at very strict $t$, $\sqrt{t}$-Greedy + DGF gives an
order-of-magnitude speedup in both regimes, and pure DGF is the
fallback that matches them on quality at the highest computational
cost.

\section{Related Work}
\label{sec:related-work}

The Descending Greedy Filter (DGF) sits at the intersection of two
lines of work on geometric spanners: \emph{(i)} constructive algorithms
that build small or light $t$-spanners directly (greedy, Yao,
$\Theta$-graphs, WSPD-based, etc.), and \emph{(ii)} structural properties
characterizing when a $t$-spanner is sparse and/or light (the
$t$-leapfrog property, existential optimality of greedy). DGF is
closest in spirit to the classical greedy spanner---it produces an
edge-minimal $t$-spanner by inspecting one edge at a time---but the
\emph{direction} of processing is reversed (descending instead of
ascending), and the algorithm is naturally a \emph{filter}: it can be
applied to any $t$-spanner, not only the complete graph. A
comprehensive reference for the whole area is the monograph of
Narasimhan and Smid~\cite{NS07}. After reviewing the greedy spanner
and its structural theory
(Section~\ref{subsec:rw-greedy}), the cone-based and Delaunay families
(Section~\ref{subsec:rw-cone}), WSPD and light spanners
(Section~\ref{subsec:rw-wspd}), faster implementations of greedy
(Section~\ref{subsec:rw-faster}), and edge-minimality, filtering, and
empirical comparisons (Section~\ref{subsec:rw-minimality}), we close
with a short summary of how DGF diverges from each of these lines
(Section~\ref{subsec:rw-we-differ}).

\subsection{Greedy spanners and the leapfrog property}
\label{subsec:rw-greedy}

The greedy spanner was introduced by Alth\"ofer, Das, Dobkin, Joseph
and Soares~\cite{ADDJS93}, who sort the edges of a weighted graph in
non-decreasing order and add an edge $(u,v)$ to the spanner only when
its current shortest-path distance exceeds $t\cdot|uv|$. They proved
that for any $k\ge 1$ the resulting $(2k-1)$-spanner has at most
$n^{1+1/k}$ edges, matching the Erd\H{o}s girth conjecture. DGF is the
natural reverse: instead of starting from the empty graph and inserting
short edges first, DGF starts from a $t$-spanner and removes long edges
first whenever the remaining subgraph is still a $t$-spanner. Both
algorithms terminate at an edge-minimal $t$-spanner, but on the same
input they generally produce \emph{different} edge-minimal subgraphs.

The structural side of greedy was developed in two important
follow-ups. Chandra, Das, Narasimhan and Soares~\cite{CDNS95} gave the
first sparseness/weight bounds for greedy in Euclidean space, showing
that the greedy $(1+\varepsilon)$-spanner has $O(n)$ edges and weight
$O((1/\varepsilon)^d\cdot \wt(\MST))$. Das, Heffernan
and Narasimhan~\cite{DHN93} introduced the now-standard
\emph{$t$-leapfrog property} and proved that any edge set satisfying it
has weight $O(\wt(\MST))$ in 3D Euclidean space; Das,
Narasimhan and Salowe~\cite{DNS95} extended this to all fixed
dimensions. The greedy spanner enjoys the leapfrog property by
construction. DGF is interesting precisely because its outputs are also
edge-minimal but do \emph{not} automatically satisfy the leapfrog
property; whether they are nevertheless light is the subject of
Conjectures~\ref{conj:dgf-weight} and~\ref{conj:min-weight}. Filtser
and Solomon~\cite{FS20} later proved that greedy is
\emph{existentially} (near-)optimal for very general input
families---its worst-case size and weight match those of \emph{any}
$(1+\varepsilon)$-spanner over the same family, up to lower-order
factors. DGF can be viewed as an alternative edge-minimal construction
whose existential and worst-case properties are open, and our empirical
comparison is in part an attempt to see how close DGF gets to the
greedy spanner's empirically optimal behavior.

\subsection{Cone-based and Delaunay spanners}
\label{subsec:rw-cone}

The cone-based family begins with Yao~\cite{Yao82}, who partitioned the
space around each point into $k$ equal-angle cones and connected each
point to its nearest neighbor in each cone---the modern \emph{Yao graph
$Y_k$}. The $\Theta$-graph variant, in which each point is connected to
the point minimizing the \emph{projected} distance onto the cone
bisector, was introduced by Clarkson~\cite{Clarkson87} and
independently by Keil~\cite{Keil88}. Keil and Gutwin~\cite{KG92} gave
the first clean $t$-spanner analysis for both, and showed that the
Delaunay triangulation itself has stretch at most
$(2\pi/3)/\cos(\pi/6)\approx 2.42$; the Delaunay triangulation was
already known to be a $t$-spanner since Dobkin, Friedman and
Supowit~\cite{DFS90}. All of these are \emph{deterministic, oblivious}
constructions: each edge is added based on local cone (or local
Delaunay) information, with no global check of shortest paths.
Consequently, Yao/$\Theta$/Delaunay graphs are usually \emph{not}
edge-minimal---many edges can be removed without violating the
$t$-spanner property---which is exactly why running DGF on top of a Yao
graph (as in our experiments) tends to drop edges aggressively.

\subsection{WSPD and light spanners}
\label{subsec:rw-wspd}

The Well-Separated Pair Decomposition (WSPD) was introduced by Callahan
and Kosaraju~\cite{CK95}, who showed that for any $s>0$ the $\Theta(n^2)$
pairs of an $n$-point set in $\mathbb{R}^d$ can be covered by
$O(s^d\cdot n)$ well-separated pairs, computable in $O(n\log n)$ time.
The WSPD immediately yields a $(1+\varepsilon)$-spanner with
$O(n/\varepsilon^d)$ edges by picking one representative edge per
well-separated pair; however, the canonical WSPD spanner has weight
$\Omega(\wt(\MST)\cdot\log n)$ unless additional pruning
is applied, and it is not edge-minimal.

A separate line of work on \emph{light} spanners is concerned with
total weight rather than edge count; a consolidated exposition of
WSPD, the leapfrog property, and light-spanner constructions is given
in Chapter~14 of Narasimhan and Smid~\cite{NS07}. Levcopoulos and
Lingas~\cite{LL92} gave the first light planar $t$-spanner of weight
$O((1+1/\varepsilon)\cdot \wt(\MST))$. The first light
$(1+\varepsilon)$-spanner in arbitrary fixed dimension was given by
Das, Narasimhan and Salowe~\cite{DNS95} via the leapfrog property,
followed by a near-quadratic implementation in Das and
Narasimhan~\cite{DN97}. Borradaile, Le and Wulff-Nilsen~\cite{BLW19}
extended the optimality of greedy to doubling metrics, and Le and
Solomon~\cite{LS25} proved tight upper and lower bounds on the size
and lightness of Euclidean $(1+\varepsilon)$-spanners, showing that
the greedy spanner achieves them. The DGF conjectures, if proved,
would place DGF outputs in the same asymptotic regime as the greedy
spanner.

\subsection{Faster greedy and approximate-greedy spanners}
\label{subsec:rw-faster}

Because the textbook greedy spanner runs in $O(n^3 \log n)$ time
na\"ively, a substantial line of work has focused on faster
implementations. Gudmundsson, Levcopoulos and Narasimhan~\cite{GLN02}
gave the first $O(n\log n)$-time algorithm producing an
\emph{approximate-greedy} spanner with $O(n)$ edges, constant degree,
and weight $O(\wt(\MST))$. The exact greedy spanner was
sped up by Bose, Carmi, Farshi, Maheshwari and Smid~\cite{BCFMS10} to
$O(n^2\log n)$ time using $O(n^2)$ space. Alewijnse, Bouts, ten~Brink
and Buchin~\cite{ABBB15} gave the first $O(n)$-space algorithm with
$O(n^2\log^2 n)$ time, computing the greedy spanner on point sets with
up to a million vertices. Our recursive binary-search variant of DGF
(Section~\ref{sec:binary-search}) is in the same spirit as these
implementation-engineering results: it batches edges and uses a
logarithmic number of APSP ``rounds'' to certify many edges at once.

\subsection{Edge-minimality, filtering, and empirical comparisons}
\label{subsec:rw-minimality}

We are not aware of any prior work that proposes the descending-order
filter algorithm we study here. DGF can be viewed as the natural
spanner-analogue of the reverse-delete algorithm for spanning trees,
with ``connected'' replaced by ``$t$-spanner of the input.'' The
complexity of computing the smallest $t$-spanner is well known to be
NP-hard: Cai~\cite{Cai94} showed that finding a minimum-cardinality
$t$-spanner is NP-complete for every fixed $t\ge 2$ even on
bounded-degree graphs, and Kortsarz~\cite{Kortsarz01} strengthened
this to polylogarithmic inapproximability: for every fixed $t\ge 5$
and every $\delta>0$, approximating the minimum-cardinality $t$-spanner
within factor $(\log n)^{1-\delta}$ is NP-hard. Sigurd and
Zachariasen~\cite{SZ04} gave an ILP-based exact algorithm for the
\emph{minimum-weight} $t$-spanner problem and observed that on
moderate-sized instances the greedy spanner is consistently within a
few percent of the optimum total weight. It is important to note
that minimum-weight and edge-minimal are \emph{different} objectives:
every minimum-weight $t$-spanner is edge-minimal, but the converse
fails, and a minimum-cardinality $t$-spanner need not be edge-minimal
either. DGF targets the polynomial-time edge-minimal objective, which
lies strictly between ``any $t$-spanner'' and the NP-hard
minimum-cardinality / minimum-weight objectives.

The notion of edge-minimality itself is implicit in the
leapfrog literature~\cite{DHN93,DNS95}: every leapfrog-property edge
set is edge-minimal in a strong sense, but the converse is not true.
From this perspective, DGF is a natural stress test of the conjecture
that minimality alone (without leapfrog) suffices for linear edge
count and $O(\wt(\MST))$ weight.

On the empirical side, Farshi and Gudmundsson~\cite{FG09} carried out
the first systematic comparison of the standard constructions (greedy,
Yao, $\Theta$, ordered $\Theta$, sink-spanner, WSPD, etc.) on random and
clustered point sets in the plane, concluding that the greedy spanner
dominates everything else on every quality metric except construction
time. Alewijnse et al.~\cite{ABBB15} extended these experiments to
point sets of up to $10^6$ vertices, confirming the same pattern. Our
experiments fit into this tradition, with the new axis being not
``constructive'' comparisons but the question of how aggressively a
filter post-pass can prune existing spanners.

\subsection{Summary: how DGF differs from the closest prior work}
\label{subsec:rw-we-differ}

Against the backdrop above, the contributions of this paper can be
located against four closely related lines in a single glance:

\begin{enumerate}
    \item \emph{Vs.\ the classical greedy spanner}~\cite{ADDJS93,CDNS95}:
        the greedy algorithm processes edges in ascending order of
        length; DGF processes them in descending order. Both produce
        edge-minimal $t$-spanners, but on the same input they in
        general output different subgraphs. Our experiments show that
        DGF-based outputs have edge counts matching greedy to within
        a fraction of a percent but total weight that is consistently
        $4$--$7\%$ lower on random Euclidean instances.
    \item \emph{Vs.\ faster greedy implementations}~\cite{BCFMS10,ABBB15}:
        those papers compute the \emph{specific} greedy spanner
        faster; DGF and its filtering variant produce an
        \emph{arbitrary} edge-minimal $t$-spanner, with the recursive
        binary-search implementation (Section~\ref{sec:binary-search})
        reducing $\APSP$ calls from $\Theta(m)$ to $O(k\log m)$.
    \item \emph{Vs.\ the minimum-weight $t$-spanner line}~\cite{SZ04,Cai94,Kortsarz01}:
        those works target $\mathsf{NP}$-hard optimization problems
        (minimum cardinality or minimum weight) and either use
        exponential-time ILPs~\cite{SZ04} or prove
        inapproximability~\cite{Cai94,Kortsarz01}. Edge-minimality is
        a strictly weaker, polynomial-time attainable objective that
        still rules out ``wasteful'' edges.
    \item \emph{Vs.\ existential-optimality results}~\cite{FS20,LS25}:
        these prove that \emph{the} greedy spanner is (near-)optimal
        in a worst-case sense. Our conjectures
        (Conjectures~\ref{conj:min-edges} and~\ref{conj:min-weight}),
        if proved, would extend the same asymptotic guarantees to
        \emph{every} edge-minimal $t$-spanner, of which the greedy
        spanner and the DGF spanners are two particular instances.
\end{enumerate}

\section{Conclusion and Open Problems}
\label{sec:open-problems}

We introduced the Descending Greedy Filter, a construction dual to the
classical greedy spanner, together with a general filtering procedure
that turns any $t$-spanner into an edge-minimal one. We gave a
recursive binary-search implementation that reduces the number of
$\APSP$ calls from $\Theta(m)$ to $O(k\log m)$. Empirically, on random
point sets with $n\in\{1000,\,2000,\,5000\}$ and
$t\in\{1.001,\,1.01,\,1.1,\,1.5\}$, DGF and its variants produce
spanners whose edge counts match the greedy spanner to within a
fraction of a percent while using $4$--$7\%$ less total weight. We
further showed that the choice of preprocessing affects only runtime
and not output quality: running DGF on top of a Yao graph is a large
practical speedup for moderate $t$ but loses its advantage at very
strict $t$ (because Yao itself ceases to be sparse), while the
$\sqrt{t}$-Greedy preprocessor remains effective in both regimes.

The main questions left open by this work are the four conjectures of
Section~\ref{sec:conjectures}. We repeat them here in order of
decreasing confidence:

\begin{enumerate}
  \item \textbf{Sparsity of DGF}
        (Conjecture~\ref{conj:dgf-edges}). Does the DGF output on every
        $n$-point planar input have $O(n)$ edges? Our experiments show
        the number of edges is consistently below $2n$ for the tested
        parameters.
  \item \textbf{Sparsity of every minimal $t$-spanner in the plane}
        (Conjecture~\ref{conj:min-edges}). Does \emph{every} edge-minimal
        planar $t$-spanner have $O(n)$ edges? This is the strongest
        form of the question and, if true, would extend Filtser and
        Solomon's~\cite{FS20} existential-optimality result from the
        greedy spanner to the entire class of edge-minimal spanners.
  \item \textbf{Weight of DGF}
        (Conjecture~\ref{conj:dgf-weight}). Is
        $\wt(E_f) = O(\wt(\MST(P)))$ for the DGF output? The
        descending-order processing appears to enforce a
        leapfrog-like property on the output, but we do not have a
        proof.
  \item \textbf{Weight of every minimal $t$-spanner}
        (Conjecture~\ref{conj:min-weight}). This is the weakest of the
        four and may in fact be false: the classical weight bound for
        greedy~\cite{CDNS95,DHN93,DNS95} rests on the $t$-leapfrog
        property, which is not implied by edge-minimality alone. A
        counterexample---a minimal $t$-spanner on $n$ points with
        weight $\omega(\wt(\MST))$---would be as interesting as a
        proof.
\end{enumerate}

A further concrete direction is to replace the $O(k\log m)$ $\APSP$ calls
of the binary-search implementation by a genuinely near-linear-time
algorithm (analogous to~\cite{GLN02,BCFMS10,ABBB15} in the greedy
setting) that incrementally maintains the shortest-path structure as
edges are removed. Finally, it would be interesting to combine DGF
with a non-trivial \emph{starting} spanner (a $\Theta$-graph or a WSPD
spanner) and study the distribution of edge-minimal outputs obtained
by different starting points.

% Bibliography is generated by BibTeX from references.bib.
% Choose the .bst that matches the target venue; see bibliography_guide.md.
%   plainurl   -- LIPIcs venues (SoCG, ESA, ICALP, ITCS, ...). Ships with
%                 the Dagstuhl LIPIcs class; not in stock TeX Live.
%   siamplain  -- SIAM venues (SODA, SICOMP, SIDMA). Ships with siamart.
%   plain      -- builtin fallback for local draft compilation.
\bibliographystyle{plainurl}
% \bibliographystyle{siamplain}
% \bibliographystyle{plain}
\bibliography{references}

\end{document}
